\begin{abstract}
	The equivalent-layer technique has been widely applied in geophysics since the 1960s 
	for potential-field transformations as well as for processing and interpreting 
	gravity and magnetic data. Despite its success, the method has traditionally relied 
	on empirical assumptions, with limited attention to its theoretical foundation. In a 
	companion paper, a rigorous theoretical framework for planar equivalent layers was 
	established. The implications of these theoretical results for the empirical 
	approaches commonly used in the geophysical literature are examined. Through 
	numerical experiments, the equivalence between the closed-form solutions for planar 
	equivalent layers and the discrete layers typically obtained by linear inversion of 
	potential-field data is demonstrated. Results confirm that multiple planar equivalent 
	layers can be employed to achieve the same potential-field transformations. An 
	application of an alternative equivalent-layer formulation is presented for airborne 
	magnetic data acquired over a complex alkaline province in central Brazil.
\end{abstract}